\documentclass[aspectratio=169, UTF8, fontset=none]{ctexbeamer}
\setCJKmainfont{FZNewShuSong-Z10}[BoldFont={Source Han Serif SC SemiBold}, ItalicFont=FZKai-Z03]
\setCJKmonofont{FZFangSong-Z02}
\newCJKfontfamily\songti{FZNewShuSong-Z10}[BoldFont={Source Han Serif SC SemiBold}]
\newCJKfontfamily\fangsong{FZFangSong-Z02}
\newCJKfontfamily\kaiti{FZKai-Z03}
\usefonttheme{serif}
\setbeamertemplate{frametitle}[default][center]
\setbeamertemplate{frametitle}{
  % \vspace{-2em}
  \bfseries\insertframetitle
  \par\smallskip
  \fangsong\large\insertframesubtitle
  \vskip-6pt\hrulefill
}
\setbeamertemplate{navigation symbols}{}
\setbeamerfont{framesubtitle}{size=\large}
\setbeamersize{text margin left=0.5em, text margin right=0.5em}
\usepackage{pgfpages}
\addtolength{\voffset}{-2em}
% \pgfpagesuselayout{resize to}[a4paper, landscape, border shrink=5mm]
\pgfpagesuselayout{4 on 1}[a4paper, landscape, border shrink=5mm]
\pgfpageslogicalpageoptions{1}{border code=\pgfusepath{stroke}}
\pgfpageslogicalpageoptions{2}{border code=\pgfusepath{stroke}}
\pgfpageslogicalpageoptions{3}{border code=\pgfusepath{stroke}}
\pgfpageslogicalpageoptions{4}{border code=\pgfusepath{stroke}}

\newcommand\hzi[1]{\textcolor{magenta}{#1}}
\newcommand\ft[1]{\smallskip\hzi{\kaiti{【法条索引】#1}}\bigskip} % 法条

\begin{document}
\begin{frame}{【1】刑事诉讼主体的范围}{3大4小2种人}
  \begin{enumerate}
    \item 3大：公、检、法；
    \item 4小：国安机关、监狱、军队保卫部门、中国海警局；
    \item 2种人：当事人、其他诉讼参与人。
  \end{enumerate}
\end{frame}

\begin{frame}{【2-1】其他诉讼参与人范围}{法、诉、辩、证、翻、鉴}
  法定代理人、诉讼代理人、辩护人、证人、翻译、鉴定人。

  \ft{（《刑诉法》第108条）}
\end{frame}

\begin{frame}{【2-2】法定代理人和诉讼代理人的区别}{法代，权限似本人，地位很独立\\诉代，权限看授权，地位只代言}
  \begin{enumerate}
    \item 法定代理人：除了个别权利（陈述权）外，可以行使本人的权利，无需被代理人本人同意。
    \item 诉讼代理人：在授权范围内，经本人同意，可以行使相应的权利，仅是被代理人的代言人。
  \end{enumerate}
\end{frame}

\begin{frame}{【3】法定不追究刑事责任的情形}{显著轻、过时效、特赦、告诉和死掉}
  \begin{enumerate}
    \item 情节显著轻微、危害不大，不认为是犯罪的；
    \item 犯罪已过追诉时效期限的；
    \item 经特赦令免除刑罚的；
    \item 依照刑法告诉才处理的犯罪，没有告诉或者撤回告诉的；
    \item 嫌疑人/被告人死亡的；
    \item 其他法律规定免予追究刑事责任的。
  \end{enumerate}

  \ft{（《刑诉法》第16条）}

  遇到这些情形的处理：
  \begin{enumerate}
    \item 立案阶段：不立案决定。
    \item 侦查阶段：撤销案件决定。
    \item 审查起诉：不起诉决定。
    \item 审判阶段：（显著轻或死者无罪）判决宣告无罪/（其他情形）裁定终止审理。
  \end{enumerate}

  \ft{（《刑诉法》第16条）}
\end{frame}

\begin{frame}{【4】作证人员的安全保障}{危恐黑毒要保护，信息音容接住他}
  【关联法条】《刑事诉讼法》第64条规定，对于危害国家安全犯罪、恐怖活动犯罪、黑社会性质的组织犯罪、毒品犯罪等案件，证人、鉴定人、被害人因在诉讼中作证，本人或者其近亲属的人身安全面临危险的，法院、检察院和公安应当采取以下一项或者多项保护措施：
  \begin{enumerate}
    \item 不公开真实姓名、住址和工作单位等个人信息；
    \item 采取不暴露外貌、真实声音等出庭作证措施；
    \item 禁止特定的人员接触证人、鉴定人、被害人及其近亲属；
    \item 对人身和住宅采取专门性保护措施；
    \item 其他必要的保护措施。
  \end{enumerate}
  证人、鉴定人、被害人认为因在诉讼中作证，本人或者其近亲属的人身安全面临危险的，可以向法院、检察院、公安请求予以保护。法院、检察院、公安依法采取保护措施，有关单位和个人应当配合。

  \ft{（《刑诉法》第64条）}
\end{frame}

\begin{frame}{【5】近亲属的范围}{上下左右}
  \begin{enumerate}
    \item 范围：父母、子女、配偶、同胞兄弟姐妹。
    \item 权限：可申请取保候审、可申请羁押必要性审查、可再审申诉、有非独立的上诉权、可为被害人委托代理人、为被告人委托辩护人等权限。
  \end{enumerate}

  \ft{（《刑诉法》第108条）}
\end{frame}

\begin{frame}{【6】同时满足下列条件证人应当出庭作证}{有异议、有影响、有必要}
  \begin{enumerate}
    \item 公诉人、当事人、辩护人、代理人对证人证言有异议；
    \item 该证言对案件定罪量刑有重大影响；
    \item 法院认为有必要。（法院通知证人出庭即视为有必要）
  \end{enumerate}

  \ft{（《刑事诉讼法》第192条）}
\end{frame}

\begin{frame}{【7】证人出庭的例外情形}{小孩、有病、太远、在国外}
  \begin{enumerate}
    \item 未成年人；
    \item 身患重病或行动不便；
    \item 远离开庭地点且交通极为不便的；
    \item 身处国外短期无法回国的；
    \item 有其他客观原因，确实无法出庭的。
  \end{enumerate}

  【考点提示】上述情形，可以通过视频方式作证。

  \ft{（《刑诉解释》第253、558条）}
\end{frame}

\begin{frame}{【8】同时满足下列条件鉴定人应当出庭作证}{有异议、有必要}
  公诉人、当事人或者辩护人、诉讼代理人对鉴定意见有异议，申请法庭通知证人、鉴定人出庭作证，法院认为有必要的。（区别证人）

  \ft{（《刑诉法》第192条）}
\end{frame}

\begin{frame}{【9-1】证人拒绝出庭的后果}{证人拒绝不出庭，强制训诫或拘留，排除拒绝不真实}
  \begin{enumerate}
    \item 经法院通知，证人没有正当理由不出庭作证的，法院可以强制其到庭，但是被告人的配偶、父母、子女除外。
    \item 证人没有正当理由拒绝出庭或者出庭后拒绝作证的，予以训诫，情节严重的，经院长批准，处以十日以下的拘留。
    \item 经法院通知，证人没有正当理由拒绝出庭或者出庭后拒绝作证，法庭对其证言的真实性无法确认的，该证人证言不得作为定案的根据。
  \end{enumerate}

  \ft{（《刑诉法》第193条《刑诉解释》第91条）}
\end{frame}

\begin{frame}{【9-2】鉴定人拒绝出庭的后果}{鉴人拒绝就排除}
  拒不出庭作证，鉴定意见不得作为定案根据。

  \ft{（《刑诉解释》第99条）}
\end{frame}

\begin{frame}{【10】监委调查的案件范围}{贪、玩、滥、舞、责、其他}
  \begin{enumerate}
    \item 贪污贿赂犯罪；
    \item 玩忽职守犯罪；
    \item 滥用职权犯罪；
    \item 徇私舞弊犯罪；
    \item 重大责任事故犯罪；
    \item 公职人员其他犯罪。（共有 6 类 101 个罪名，回关键词“101”至公众号“高甲伴你过法考”可查询）
  \end{enumerate}

  \ft{（《监察法实施条例》第26-31条）}
\end{frame}

\begin{frame}{【11】检察院自侦案件的第一类}{搜、拘、滥、私、逼；暴、虐、枉、玩、执；放、假监、脱逃}
  检察院在对诉讼活动实行法律监督中发现的司法工作人员利用职权实施的非法拘禁、刑讯逼供、非法搜查等侵犯公民权利、损害司法公正的犯罪，可以由检察院立案侦查。
  【具体罪名】
  \begin{enumerate}
    \item 非法搜查罪（非司法工作人员除外）；
    \item 非法拘禁罪（非司法工作人员除外）；
    \item 滥用职权罪（非司法工作人员滥用职权除外）；
    \item 徇私枉法罪；
    \item 刑讯逼供罪；
    \item 暴力取证罪；
    \item 虐待被监管人罪；
    \item 民事、行政枉法裁判罪；
    \item 玩忽职守罪（非司法工作人员玩忽职守除外）；
    \item 执行判决、裁定失职罪；
    \item 执行判决、裁定滥用职权罪；
    \item 私放在押人员罪；
    \item 徇私舞弊减刑、假释、暂予监外执行罪；
    \item 失职致使在押人员脱逃罪。
  \end{enumerate}

  \ft{（《刑诉法》第19条第2款）}
\end{frame}

\begin{frame}{【12】检察院自侦案件的第二类}{公、机、权、大、省以上}
  对于公安管辖的国家机关工作人员利用职权实施的重大犯罪案件，需要由检察院直接受理的时候，经省级以上检察院决定，可以由检察院立案侦查。

  \ft{（《刑诉法》第19条第2款）}
\end{frame}

\begin{frame}{【13-1】自诉案件的全部类型}{亲告、轻微、公转自}
  \begin{enumerate}
    \item 告诉才处理的案件（亲告）
    \item 被害人有证据证明的轻微刑事案件（轻微）
    \item 公诉转自诉案件（公转自）
  \end{enumerate}

  \ft{（《刑诉解释》第1条）}
\end{frame}

\begin{frame}{【13-2】告诉才处理的案件范围}{侮、诽、暴、虐、侵}
  \begin{enumerate}
    \item 侮辱、诽谤案（刑法第246条规定的，但严重危害社会秩序和国家利益的除外）；
    \item 暴力干涉婚姻自由案（刑法第257条第1款规定的）；
    \item 虐待案（刑法第160条第1款规定的，但被害人没有能力告诉或者因受到强制、威吓无法告诉的除外）；
    \item 侵占案（刑法第270条规定的）。
  \end{enumerate}

  \ft{（《刑诉解释》第1条第1项）}
\end{frame}

\begin{frame}{【13-3】检察院没有提起公诉，被害人有证据证明的轻微刑事案件（可公可自）}{重、轻、伪、弃、通、知、住、三下}
  \begin{enumerate}
    \item 故意伤害案；
    \item 非法侵入住宅案；
    \item 侵犯通信自由案；
    \item 重婚案；
    \item 遗弃案；
    \item 生产、销售伪劣商品案；
    \item 侵犯知识产权案；
    \item 刑法分则第 4 章、第 5 章规定的，可能判处 3 年有期徒刑以下刑罚的案件。
  \end{enumerate}

  \ft{（《刑诉解释》第1条第2项）}
\end{frame}

\begin{frame}{【13-4】第二类自诉案件（可公可自）的特点}{轻微案件，可公可自，正在公诉，不能自诉}
  此类自诉案件可以公诉也可以自诉，选择公诉或者自诉没有先后顺序。但是如果公安正在立案侦查或检察院正在审查起诉的，自诉人提起自诉法院不予受理。

  \ft{（《刑诉解释》第320条第2款第8项）}
\end{frame}

\begin{frame}{【14-1】普通案件与监察机关的管辖竞合}{普监交叉，该谁给谁，监察为主}
  公、检、法各机关在办案中发现应当由其他机关侦查或调查的案件，应当移送至该机关。被调查人既涉嫌严重职务违法或者职务犯罪，又涉嫌其他违法犯罪的，一般应当由监察机关为主调查，其他机关予以协助。

  \ft{（《监察法》第34条）}
\end{frame}

\begin{frame}{【14-2】公安与检察院的管辖竞合}{公检交叉，该谁给谁，再看主罪}
  将不属于本部门管辖的案件移送给对方。在上述情况中，如果涉嫌主罪属于其中一个机关管辖，则由该机关为主侦查，另一侦查机关予以配合。

  \ft{（《高检规则》第18条或《公安部规定》第30条）}
\end{frame}

\begin{frame}{【14-3】公诉中发现有自诉交叉}{公自交叉，该谁给谁}
  \begin{enumerate}
    \item 属于告诉才处理案件，可告知被害人向法院直接自诉。
    \item 属于其他自诉案件，可直接立案进行侦查，和公诉案件一并移送法院，由法院合并审理。
  \end{enumerate}
\end{frame}

\begin{frame}{【14-4】自诉中发现有公诉交叉}{自公交叉，该谁给谁}
  法院在审理自诉案件过程中，如果发现被告人还犯有必须由检察院提起公诉的罪行时，则应将新发现的罪行另案移送有管辖权的公安或者检察院或监察委处理。
\end{frame}

\begin{frame}{【15】并案管辖的几种情形}{数罪数人存关联；共犯还有其他罪}
  具有下列情形之一的，公、检、法机关可以在其职责范围内并案处理：
  \begin{enumerate}
    \item 一人犯数罪的；共同犯罪的；
    \item 共同犯罪的嫌疑人/被告人还实施其他犯罪的；
    \item 多个嫌疑人/被告人实施的犯罪存在关联，并案处理有利于查明案件事实的。
  \end{enumerate}

  \ft{（《六机关规定》第3条、《刑诉解释》第24、25条）}
\end{frame}

\begin{frame}{【16】中院管辖范围}{危、恐、无、死、没、缺席}
  \begin{enumerate}
    \item 危害国家安全、恐怖活动案件；
    \item 可能判处无期徒刑、死刑的案件；
    \item 违法所得没收程序；
    \item 缺席审判程序。
  \end{enumerate}

  \ft{（《刑诉法》第21、291、299条）}
\end{frame}

\begin{frame}{【17】级别管辖的变通}{上可以审下，下不可审上\\数罪或数人，就高不就低}
  \begin{enumerate}
    \item 上级法院必要时，可以审判下级法院管辖的一审案件；
    \item 应当由上级法院管辖的，不能再指定下级法院管辖；
    \item 只要一人或一罪属于上级法院管辖，全案由上级法院管辖。
  \end{enumerate}

  \ft{（《刑诉解释》第15-17条）}
\end{frame}

\begin{frame}{【18】地域管辖}{犯罪地为主，居住地为辅；共同管辖，先到先得}
  \begin{enumerate}
    \item 以犯罪地法院管辖为主，被告人居住地法院管辖为辅。
    \item 共同管辖，以最初受理的法院审判为主，必要时可以在主要犯罪地法院审判。
  \end{enumerate}

  \ft{（《刑诉法》第25、26条）}
\end{frame}

\begin{frame}{【19-1】指定管辖的适用情形}{管辖不明、管辖不宜、规避管辖}
  \begin{enumerate}
    \item 地域管辖不明的刑事案件。
    \item 有管辖权的法院不宜行使管辖权的刑事案件。
    \item 检察院规避法院管辖的情形。
  \end{enumerate}

  \ft{（《刑诉解释》第18、19、23条）}
\end{frame}

\begin{frame}{【19-2】指定管辖后的案件移送问题}{公诉要退回，自诉直接送}
  \begin{enumerate}
    \item 对于公诉案件：书面通知检察院，案卷材料退回，同时书面通知当事人；
    \item 对于自诉案件：将全部案卷移送被指定管辖的法院，并书面通知当事人。
  \end{enumerate}

  \ft{（《刑诉解释》第22条）}
\end{frame}

\begin{frame}{【20-1】外国人犯我国缔结或参加的国际条约规定的罪行（普遍管辖）}{抓获、登陆、入境地}
  由被告人被抓获地、登陆地或者入境地的法院管辖。

  \ft{（《刑诉解释》第12条）}
\end{frame}

\begin{frame}{【20-2】外国人在国外对中国（公民）犯罪（保护管辖）}{被告登陆和两境；被害可有现前管}
  由该外国人登陆地、入境地或入境后居住地管辖，也可由被害人离境前居住地或现居住地管辖。

  \ft{（《刑诉解释》第11条）}
\end{frame}

\begin{frame}{【21-1】中国公民在驻外领使馆犯罪（司法豁免）}{打回原籍才能管}
  由公民主管单位所在地或原户籍所在地的法院管辖。

  \ft{（《刑诉解释》第9条）}
\end{frame}

\begin{frame}{【21-2】中国公民在领域外犯罪（属人管辖）}{被告登陆、入境地；双方都有现前管}
  由被告人登陆地、入境地、离境前居住地或现居住地的法院管；被害人是中国公民的，也可由被害人离境前居住地或现居住地的法院管。

  \ft{（《刑诉解释》第10条）}
\end{frame}

\begin{frame}{【22】在中国领域外的中国航空器内的犯罪}{飞机最初降落地}
  由该航空器在中国最初降落地的法院管辖。

  \ft{（《刑诉解释》第8条）}
\end{frame}

\begin{frame}{【23-1】在中国领域外的中国船舶内的犯罪}{初停、登陆、入境地}
  由该船舶最初停泊的中国口岸所在地或被告人登陆地、入境地的法院管辖。

  \ft{（《刑诉解释》第7条）}
\end{frame}

\begin{frame}{【23-2】在内水、领海发生的刑事案件}{普通管辖 + 登陆}
  由犯罪地或被告人登陆地法院管辖。被告人居住地更适宜的，可由被告人居住地的法院管辖。

  \ft{（《刑诉解释》第4条）}
\end{frame}

\begin{frame}{【24-1】在国际列车上的犯罪}{没有协议始停靠}
  没有协定的，由该列车始发或前方停靠的中国车站所在地法院管辖。

  \ft{（《刑诉解释》第6条）}
\end{frame}

\begin{frame}{【24-2】在国内列车上的犯罪}{运行中抓获：前方停靠、始终管；\\非运行抓获：铁路公安对应管；\\在车站抓获：也可就地管}
  \begin{enumerate}
    \item 被告人在列车运行途中被抓获的，由前方停靠站所在地的法院管辖。必要时，也可以由始发站或者终点站所在地的法院管辖。
    \item 被告人不是在列车运行途中被抓获的，由负责该列车乘务的铁路公安对应的法院管辖。
    \item 被告人在列车运行途经车站被抓获的，也可以由该车站所在地的法院管辖。
  \end{enumerate}

  \ft{（《刑诉解释》第5条）}
\end{frame}

\begin{frame}{【25】漏罪、新罪的管辖}{漏罪首选原审地；新罪首选服刑地}
  \begin{enumerate}
    \item 漏罪的管辖：原则上为原审地法院管辖；（服刑地、新发现罪的主要犯罪地更为适宜的也可管辖）
    \item 新罪的管辖：原则上为服刑地法院管辖；（脱逃期犯罪，犯罪地抓获并发现，由犯罪地法院管辖）
  \end{enumerate}

  \ft{（《刑诉解释》第13条）}
\end{frame}

\begin{frame}{【26】回避的对象}{侦、检、审、书、翻、鉴}
  侦查人员、检察人员、审判人员、书记员、翻译、鉴定人（专门知识的人参照鉴定人）

  \ft{（《刑诉法》第29、32条）}
\end{frame}

\begin{frame}{【27-1】回避的理由}{利害关系；请客送礼；参前不参后}
  \begin{enumerate}
    \item 利害关系 ：见《刑诉解释》第27条
    \item 请客送礼 ：见《刑诉解释》第28条
    \item 参前不参后：见《刑诉解释》第29条《高检规则》第35条
  \end{enumerate}
\end{frame}

\begin{frame}{【27-2】不受回避规定限制的情形}{发回重组，回来不限}
  发回重新审判的案件，在第一审法院作出裁判后又进入第二审程序、在法定刑以下判处刑罚的复核程序或者死刑复核程序的，原第二审程序、在法定刑以下判处刑罚的复核程序或者死刑复核程序中的合议庭组成人员不受回避规定的限制。

  \ft{（《刑诉解释》第29条）}
\end{frame}

\begin{frame}{【28-1】申请回避、排非的主体}{当、法、诉、辩}
  申请回避的主体：当事人、法定代理人、诉讼代理人、辩护人（无近亲属）。

  \ft{（《刑诉法》第29、32条）}

  申请排非的主体：当事人、法定代理人、诉讼代理人、辩护人（无近亲属）。

  \ft{（《刑诉法》第58条）}
\end{frame}

\begin{frame}{【28-2】申请变更强制措施}{当、法、近、辩}
  嫌疑人被告人、法定代理人、近亲属、辩护人（含值班律师）

  \ft{（《刑诉法》第97条）}
\end{frame}

\begin{frame}{【28-3】再审申诉的主体}{当、法、近、外}
  当事人、法定代理人、近亲属、案外人

  \ft{（《刑诉解释》第451条）}
\end{frame}

\begin{frame}{【29】申请回避的效果}{一般要暂停，侦查才除外}
  被申请回避的人员一般应暂停参与本案的诉讼活动；但侦查人员例外。

  \ft{（《刑诉法》第31条第2款）}
\end{frame}

\begin{frame}{【30】回避的决定主体}{一般找老大，老大找组织}
  \begin{enumerate}
    \item 审判人员、检察人员、侦查人员的回避，应当分别由院长、检察长、公安负责人决定；
    \item 院长的回避，由本院审委会决定；检察长和公安负责人的回避，由同级检察院检察委员会决定。
  \end{enumerate}

  \ft{（《刑诉法》第31条第1款）}
\end{frame}

\begin{frame}{【31】申请出庭的检察人员回避的处理方式}{有理则休庭，无理则驳回}
  当事人及其法定代理人申请出庭的检察人员回避的，法院应当区分情况作出处理：
  \begin{enumerate}
    \item 属于刑诉法规定情形的回避申请，应当决定休庭，并通知检察院尽快作出决定；
    \item 不属于刑诉法规定情形的回避申请，应当当庭驳回，并不得申请复议。
  \end{enumerate}

  \ft{（《刑诉解释》第36条）}
\end{frame}

\begin{frame}{【32】辩护人的人数}{1人最多找2辩，1人只为1人辩}
  \begin{enumerate}
    \item 嫌疑人/被告人除自己行使辩护权以外，还可以委托1至2人作为辩护人。
    \item 1名辩护人不得为2名以上同案被告人，或未同案处理但犯罪存在关联的被告人辩护。
  \end{enumerate}

  \ft{（《刑诉解释》第43条）}
\end{frame}

\begin{frame}{【33-1】辩护人的绝对禁止范围}{刑罚执行中、自由受限、能力有限}
  \begin{enumerate}
    \item 正在被执行刑罚或者处于缓刑、假释考验期间的人；（刑罚执行中）
    \item 依法被剥夺、限制人身自由的人；（自由受限制）
    \item 无行为能力或者限制行为能力的人；（能力有限制）
  \end{enumerate}

  \ft{（《刑诉解释》第40条）}
\end{frame}

\begin{frame}{【33-2】辩护人的相对禁止范围}{现、陪、利、外、吊执照}
  \begin{enumerate}
    \item 法院、检察院、公安、国安、监狱、监察机关的现职人员；
    \item 人民陪审员；
    \item 与本案审理结果有利害关系的人；
    \item 外国人或者无国籍人。
    \item 被开除公职和被吊销律师、公证员执业证书的人。
  \end{enumerate}

  【考点提示】上述情形如果是被告人的近亲属或者监护人，仍可作为其辩护人。

  \ft{（《刑诉解释》第40条）}
\end{frame}

\begin{frame}{【33-3】相对禁止担任辩护人的“现职”的范围}{公、检、法、安、监}
  法院、检察院、公安、国家安全机关、监狱、监察机关的“现职”人员一般禁止担任辩护人。

  \ft{（《刑诉解释》第40条）}
\end{frame}

\begin{frame}{【34】法官、检察官及其至亲担任辩护人、诉讼代理人的限制}{两年、终身、家庭店}
  \begin{enumerate}
    \item 两年：从法院、检察院离任后2年内，不得以律师身份担任诉讼代理人或辩护人。
    \item 终身：
          \begin{enumerate}
            \item 离任后，不得担任原任职机关办理案件的诉讼代理人或者辩护人，但是作为当事人的监护人或近亲属代理诉讼或进行辩护的除外。
            \item 被开除后，不得担任诉讼代理人或者辩护人，但是作为当事人的监护人或近亲属代理诉讼或进行辩护的除外。
          \end{enumerate}
    \item 家庭店：
          \begin{enumerate}
            \item 法官、检察官的配偶、父母、子女有下列情形之一的，应当实行任职回避：
            \item 担任其任职机关辖区内律师事务所的合伙人或者设立人的；
            \item 在其任职机关辖区内以律师身份担任诉讼代理人、辩护人，或为诉讼案件当事人提供其他有偿法律服务的。
          \end{enumerate}
  \end{enumerate}

  \ft{（《法官法》36、24条《检察官法》37、25条）}
\end{frame}

\begin{frame}{【35】非律师辩护人和律师辩护人权利的差别}{非律辩护矮半截}
  \begin{enumerate}
    \item 非律阅卷要许可，律师阅卷不用批。
    \item 非律会见要许可，律师一般不用批。
    \item 非律没有取证权，律师取证需同意。
  \end{enumerate}

  \ft{（《刑诉法》第39、40、43条）}
\end{frame}

\begin{frame}{【36】侦查阶段的消极权利清单}{侦查阶段无阅卷、无核证、无非律}
  \begin{enumerate}
    \item 侦查阶段的消极权利清单：
          \begin{enumerate}
            \item 侦查阶段没有阅卷权；
            \item 侦查阶段会见时不能核实证据；
            \item 侦查阶段不能出现非律师辩护人。
          \end{enumerate}
    \item 侦查阶段律师辩护人的权利清单：
          \begin{enumerate}
            \item 可以提供帮助；
            \item 可以代理申控；
            \item 可以申请变更；
            \item 可以了解情况；
            \item 可以会见通信。
          \end{enumerate}
  \end{enumerate}

  \ft{（《刑诉法》第38、39条）}
\end{frame}

\begin{frame}{【37】辩护人的阅卷权在不同身份和不同阶段的差别}{非律阅卷要许可，律师阅卷不用批，均自审查起诉起}
  辩护律师自检察院对案件审查起诉之日起，可以查阅、摘抄、复制本案的案卷材料。其他辩护人经法院、检察院许可，也可以查阅、摘抄、复制上述材料。

  \ft{（《刑诉法》第40条）}
\end{frame}

\begin{frame}{【38-1】辩护人的会见、通信权在不同身份上的差别}{非律会见要许可，律师一般不用批}
  辩护律师可以同在押的嫌疑人/被告人会见和通信。其他辩护人经法院、检察院许可，也可以同在押的嫌疑人/被告人会见和通信。

  \ft{（《刑诉法》第39条第1款）}
\end{frame}

\begin{frame}{【38-2】辩护律师的会见权、核实证据权在诉讼阶段上的差别}{侦查律师可会见，核证等到下阶段}
  辩护律师侦查阶段可以会见在押的或监视居住的嫌疑人/被告人，可以了解案件有关情况，提供法律咨询。自案件移送审查起诉之日起，可以向嫌疑人/被告人核实有关证据。

  \ft{（《刑诉法》第39条第4款）}
\end{frame}

\begin{frame}{【38-3】律师会见权在不同案件上的差别}{律师会见凭3证，危恐除外要许可}
  \begin{enumerate}
    \item 辩护律师持律师执业证书、律师事务所证明和委托书或者法律援助公函要求会见在押的嫌疑人/被告人的，看守所应当及时安排会见，至迟不得超过48小时。
    \item 危害国家安全犯罪、恐怖活动犯罪，在侦查期间辩护律师会见在押的或监视居住的嫌疑人，应当经侦查机关许可。
  \end{enumerate}

  \ft{（《刑诉法》第39条第2、3款）}
\end{frame}

\begin{frame}{【39】通信可被截留、复制、删改的特殊情形}{国公人身串毁证，信件可被截复删}
  看守所可以对信件进行必要的检查，但不得截留、复制、删改信件，不得向办案机关提供信件内容，但信件内容涉及危害国家安全、公共安全、严重危害他人人身安全以及涉嫌串供、毁灭证据等情形的除外。

  \ft{（《关于依法保障律师执业权利的规定》第13条）}
\end{frame}

\begin{frame}{【40】辩护人的取证权在不同身份上的差别}{非律没有取证权，律师取证需同意，控方证人双重许，若有困难申请取}
  \begin{enumerate}
    \item 辩护律师经证人或者其他有关单位和个人同意，可以向他们收集与本案有关的材料；
    \item 辩护律师经检察院或法院许可，并且经被害人或其近亲属、被害人提供的证人同意，可以向他们收集与本案有关的材料；
    \item 辩护律师也可以申请检察院、法院收集、调取证据，或者申请法院通知证人出庭作证。
  \end{enumerate}

  \ft{（《刑诉法》第43条）}
\end{frame}

\begin{frame}{【41】主动应当听取辩护人意见的情形}{批捕未成年，审查起诉中，二审不开庭，死刑上抗案，认罪又认罚。}
  \begin{enumerate}
    \item 审查批捕未成年人，必须听辩护人意见。
    \item 检察院审查起诉，应当讯问嫌疑人，听取辩护人、被害人及其诉讼代理人的意见。
    \item 二审决定不开庭审理的，应当讯问被告人，听取其他当事人、辩护人、诉讼代理人的意见。
    \item 死刑上诉、抗诉案件，检察院应当讯问原审被告人，应当听取辩护人意见。
    \item 认罪认罚案件，应当听取被告人及其辩护人或者值班律师的意见。
  \end{enumerate}

  \ft{（《刑诉法》第280、173、234条《高检规则》第450条）}
\end{frame}

\begin{frame}{【42】辩护人的证据开示义务}{不在场、不够大、不正常}
  辩护人收集的有关嫌疑人不在犯罪现场、未达到刑事责任年龄、属于依法不负刑事责任的精神病人的证据，应当及时告知公安、检察院。

  \ft{（《刑诉法》第42条）}
\end{frame}

\begin{frame}{【43】辩护律师的保密义务}{过去统统要保密，将来国、公、人身要揭发}
  辩护律师对在执业活动中知悉的委托人的有关情况和信息，有权予以保密。但是，辩护律师在执业活动中知悉委托人或者其他人，准备或者正在实施危害国家安全、公共安全以及严重危害他人人身安全的犯罪的，应当及时告知司法机关。

  \ft{（《刑诉法》第48条）}
\end{frame}

\begin{frame}{【44】强制法援辩护的情形}{盲聋哑、半疯傻、死无缺、未长大}
  \begin{enumerate}
    \item 盲、聋、哑；
    \item 尚未完全丧失辨认或者控制自己行为能力的精神病人；
    \item 可能判处无期、死刑的；缺席审判程序；
    \item 诉讼时未满18周岁的未成年人。
  \end{enumerate}

  \ft{（《刑诉解释》第47条、《刑诉法》278、293条）}
\end{frame}

\begin{frame}{【45】拒绝辩护的考点}{有理、1次、有人辩；无理、2次、自己辩}
  \begin{enumerate}
    \item 强制款：必须有理由，且只能拒绝1次，拒绝之后必须保证有人为他辩护；
    \item 普通款：不用理由，可以拒绝2次，两次拒绝之后只能自己辩护。
  \end{enumerate}

  \ft{（《刑诉解释》第50、311条）}
\end{frame}

\begin{frame}{【46】值班律师的职责}{提供咨询出建议申请变更提意见申请法援或其他}
  \begin{enumerate}
    \item 提供法律咨询；
    \item 提供程序选择建议；
    \item 帮助嫌疑人/被告人申请变更强制措施；
    \item 对案件处理提出意见；
    \item 帮助嫌疑人/被告人及其近亲属申请法律援助；
    \item 法律法规规定的其他事项。
  \end{enumerate}

  【关联考点】值班律师在阅卷、会见、意见权上同辩护律师一样，但是不能出庭辩护。

  \ft{（《关于适用认罪认罚从宽制度的指导意见》第12条）}
\end{frame}

\begin{frame}{【47】证据需要掌握的主要内容}{3-3-8-4-7}
  3大属性；3大原则；8大法定种类；4大理论分类；7大证据规则。
\end{frame}

\begin{frame}{【48】物证、书证强制排除的情形}{真伪不明、来源不明、无法解释}
  \begin{enumerate}
    \item 物证的照片、录像、复制品，不能反映原物的外形和特征的；（真伪不明）
    \item 书证有更改或者更改迹象不能作出合理解释，或书证的副本、复制件不能反映原件及其内容的；（真伪不明）
    \item 在勘验、检查、搜查过程中提取、扣押的物证、书证，未附笔录或者清单，不能证明物证、书证来源的；（来源不明）
    \item 对物证、书证的来源、收集程序有疑问，不能作出合理解释的。（无法解释）
  \end{enumerate}

  \ft{（《刑诉解释》第83-86条）}
\end{frame}

\begin{frame}{【49】证人证言强制排除的情形}{麻醉、猜测、未个别；核对、翻译、暴限胁；拒绝出庭不真实}
  \begin{enumerate}
    \item 醉酒、中毒或者麻醉等状态，不能正常感知或者正确表达的证人所提供的证言。
    \item 证人的猜测性、评论性、推断性的证言，但根据一般生活经验判断符合事实的除外。
    \item 询问证人没有个别进行的，不得作为证据使用。
    \item 书面证言没有经证人核对确认的，不得作为证据使用。
    \item 应当提供翻译人员而未提供的，不得作为证据使用。
    \item 采用暴力、非法限制人身自由、威胁以及等非法方法收集的证人证言，不得作为证据使用。
    \item 证人没有正当理由拒绝出庭或者出庭后拒绝作证，法庭对其证言的真实性无法确认。
  \end{enumerate}

  【考点提示】询问未成年人，其法定代理人或合适成年人不在场的，属于可以补正、解释的情形。

  \ft{（《刑诉解释》第88、89、91、125条）}
\end{frame}

\begin{frame}{【50】供述强制排除的情形}{核对、翻译、无法代；场外、音像、暴限胁}
  \begin{enumerate}
    \item 讯问笔录没有经被告人核对确认的；
    \item 应当提供通晓聋、哑手势的人员或翻译人员而未提供的；
    \item 讯问未成年人，其法定代理人或合适成年人不在场的；
    \item 除情况紧急必须现场讯问以外，在规定的办案场所外讯问取得的供述；
    \item 未依法对讯问进行全程录音录像取得的供述；
    \item 采用暴力方法或变相肉刑手段，使被告人遭受难以忍受的痛苦而违背意愿作出的供述；
    \item 采用非法限制人身自由的方法收集的供述。
    \item 采用以暴力或严重损害本人及其近亲属合法权益等相威胁的方法，使被告人遭受难以忍受的痛苦而违背意愿作出的供述；
  \end{enumerate}

  【考点提示】讯问未成年人，其法定代理人或合适成年人不在场的，属于强制排除的情形。

  \ft{（《刑诉解释》第94、123条《防范冤假错案工作机制意见》第8条）}
\end{frame}

\begin{frame}{【51】辨认笔录强制排除的情形}{不是侦调来主持；指示预见个混混}
  \begin{enumerate}
    \item 辨认不是在调查人员、侦查人员主持下进行的；
    \item 辨认前使辨认人见到辨认对象的；
    \item 辨认活动没有个别进行的；
    \item 辨认对象没有混杂在同类对象中，或供辨认的对象数量不符合规定的；（公7人5物10照）
    \item （检7人5物10人照5物照）
    \item 辨认中给辨认人明显暗示或者明显有指认嫌疑的；
    \item 其他违反规定、不能确定辨认笔录真实性的。
  \end{enumerate}

  \ft{（《刑诉解释》第105条）}
\end{frame}

\begin{frame}{【52】鉴定意见强制排除的情形：（略）}{见错就排}
  \ft{（《刑诉解释》第98条）}
\end{frame}

\begin{frame}{【53】勘验、检查笔录强制排除的排除}{无法解释}
  存在明显不符合法律、有关规定的情形，不能作出合理解释或说明的。

  \ft{（《刑诉解释》第103条）}
\end{frame}

\begin{frame}{【54】侦查实验笔录强制排除的情形}{条件差异}
  侦查实验的条件与事件发生时的条件有明显差异，或存在影响实验结论科学性的其他情形的。

  \ft{（《刑诉解释》第107条）}
\end{frame}

\begin{frame}{【55】视听资料强制排除的情形}{真伪不明、无法解释}
  视听资料强制排除的情形：
  \begin{enumerate}
    \item 经审查无法确定真伪的；
    \item 制作、取得的时间、地点、方式等有疑问，不能提供必要证明或者作出合理解释的。
  \end{enumerate}

  \ft{（《刑诉解释》第109条）}

  电子数据强制排除的情形：
  \begin{enumerate}
    \item 系篡改、伪造或者无法确定真伪的；
    \item 有增加、删除、修改等情形，影响电子数据真实性的；
    \item 其他无法保证电子数据真实性的情形。
    \item 若有瑕疵，不能补正或者作出合理解释的，不得作为定案的根据。
  \end{enumerate}

  \ft{（《刑诉解释》第113、114条）}
\end{frame}

\begin{frame}{【56】非法实物可补正}{非法实物可补正，言词排非暴限胁，重复供述一并，排除非换人换阶段}
  \begin{enumerate}
    \item 非法实物可补正。物证、书证不符合法定程序，可能严重影响司法公正的，应当予以补正或作出合理解释；不能补正或解释的，应当予以排除。
    \item 言词排非暴限胁：
          \begin{enumerate}
            \item 《刑诉解释》第123条 采用下列非法方法收集的被告人供述，应当予以排除：
                  \begin{enumerate}
                    \item 采用殴打、违法使用戒具等暴力方法或者变相肉刑的恶劣手段，使被告人遭受难以忍受的痛苦而违背意愿作出的供述；
                    \item 采用以暴力或者严重损害本人及其近亲属合法权益等相威胁的方法，使被告人遭受难以忍受的痛苦而违背意愿作出的供述；
                    \item 采用非法拘禁等非法限制人身自由的方法收集的被告人供述。
                  \end{enumerate}
            \item 《刑诉解释》第125条 采用暴力、威胁以及非法限制人身自由等非法方法收集的证人证言、被害人陈述，应当予以排除。
          \end{enumerate}
    \item 重复供述一并排，除非换人换阶段：《刑诉解释》第124条 采用刑讯逼供方法使被告人作出供述，之后被告人受该刑讯逼供行为影响而作出的与该供述相同的重复性供述，应当一并排除，但下列情形除外：
          \begin{enumerate}
            \item 调查、侦查期间，监察机关、侦查机关根据控告、举报或者自己发现等，确认或者不能排除以非法方法收集证据而更换调查、侦查人员，其他调查、侦查人员再次讯问时告知有关权利和认罪的法律后果，被告人自愿供述的；
            \item 审查逮捕、审查起诉和审判期间，检察人员、审判人员讯问时告知诉讼权利和认罪的法律后果，被告人自愿供述的。
          \end{enumerate}
  \end{enumerate}

  \ft{（《刑诉法》第56条《刑诉解释》第123-125条）}
\end{frame}

\begin{frame}{【57】免证对象的范围}{常识、法规、和定律，推定、生效、无异议}
  \begin{enumerate}
    \item 为一般人共同知晓的常识性事实；
    \item 法院生效裁判所确认的并且未依审判监督程序重新审理的事实；
    \item 法律、法规的内容以及适用等属于审判人员履行职务所应当知晓的事实；
    \item 在法庭审理中不存在异议的程序事实；
    \item 法律规定的推定事实；
    \item 自然规律或者定律。
  \end{enumerate}

  \ft{（《高检规则》第401条）}
\end{frame}

\begin{frame}{【58-1】可以取保候审的情形}{徒刑以下、不危险；疾病孕乳、超期限}
  \begin{enumerate}
    \item 可能判处管制、拘役或者独立适用附加刑的；
    \item 可能判处有期徒刑以上，采取取保候审不致发生社会危险性的；
    \item 严重疾病、生活不能自理，怀孕或正在哺乳自己婴儿，不致发生社会危险性的；
    \item 羁押期限届满，案件尚未办结，需要采取取保候审的；
    \item 其他需变更逮捕措施的。
  \end{enumerate}

  \ft{（《刑诉法》第67条）}
\end{frame}

\begin{frame}{【58-2】禁止取保候审的情形}{禁止： 累、主、残、暴、严；\\除非：疾病、孕乳、超期限}
  对累犯，犯罪集团的主犯，以自伤、自残办法逃避侦查的嫌疑人，严重暴力犯罪以及其他严重犯罪的嫌疑人不得取保候审，但嫌疑人具有严重疾病、生活不能自理，怀孕或正在哺乳自己婴儿，不致发生社会危险性的或者羁押期限届满，案件尚未办结，需要采取取保候审的情形除外。

  \ft{（《公安部规定》第82条）}
\end{frame}

\begin{frame}{【59】适用保证人取保候审的情形}{老、幼、没钱}
  \begin{enumerate}
    \item 无力交纳保证金的；
    \item 未成年人或已满75周岁的人；
    \item 其他不宜收取保证金的情形。
  \end{enumerate}

  \ft{（《刑诉解释》第151条）}
\end{frame}

\begin{frame}{【60】保证人的基本条件}{无牵连、有能力、有权、有房又有钱}
  \begin{enumerate}
    \item 无牵连；
    \item 有能力履行保证义务；
    \item 享有政治权利，自由未受到限制；
    \item 有固定的住处和收入。
  \end{enumerate}

  \ft{（《刑诉法》第69条）}
\end{frame}

\begin{frame}{【61】保证人的义务和责任}{监督、报告、罚款、刑责}
  \begin{enumerate}
    \item 监督被保证人遵守刑诉法规定的义务；
    \item 发现被保证人可能发生或已经发生违反法定行为的，应当及时向执行机关报告；
    \item 被保证人有违反刑诉法规定的行为，保证人未履行保证义务的，对保证人处以罚款；
    \item 构成犯罪的，依法追究刑事责任。
  \end{enumerate}

  \ft{（《刑诉法》第70条）}
\end{frame}

\begin{frame}{【62-1】被取保候审的嫌疑人/被告人应当遵守的5项义务}{市县、报告、及时到、不扰证人、不毁证}
  \begin{enumerate}
    \item 未经执行机关批准不得离开所居住的市、县；
    \item 住址、工作单位和联系方式发生变动的，在24小时以内向执行机关报告；
    \item 在传讯的时候及时到案；
    \item 不得以任何形式干扰证人作证；
    \item 不得毁灭、伪造证据或者串供。
  \end{enumerate}

  \ft{（《刑诉法》第71条第1款）}
\end{frame}

\begin{frame}{【62-2】可以要求被取保候审的嫌疑人/被告人遵守的4项义务}{特定地方、特定人、特定活动、特定证}
  \begin{enumerate}
    \item 不得进入特定的场所；
    \item 不得与特定的人员会见或者通信；
    \item 不得从事特定的活动；
    \item 将护照等出入境证件、驾驶证件交执行机关保存。
  \end{enumerate}

  \ft{（《刑诉法》第71条第2款）}
\end{frame}

\begin{frame}{【63】可用监视居住替代逮捕的情形}{替捕：疾病、孕乳、唯一扶，特殊、需要、已超期\\替保：没钱又没人}
  \begin{enumerate}
    \item 可用监视居住替代逮捕的情形：公检法对符合逮捕条件，有下列情形之一的嫌疑人/被告人，可以监视居住：
          \begin{enumerate}
            \item 患有严重疾病、生活不能自理的；
            \item 怀孕或者正在哺乳自己婴儿的妇女；
            \item 系生活不能自理的人的唯一扶养人；
            \item 因为案件的特殊情况或者办理案件的需要，采取监视居住措施更为适宜的；
            \item 羁押期限届满，案件尚未办结，需要采取监视居住措施的。
          \end{enumerate}
    \item 可用监视居住替代取保候审的情形：对符合取保候审条件，但嫌疑人/被告人不能提出保证人，也不交纳保证金的，可以监视居住。
  \end{enumerate}

  \ft{（《刑诉法》第74条）}
\end{frame}

\begin{frame}{【64】被监视居住人应当遵守的6项义务}{处所、会见、及时到、不扰证人、不毁证、保存还要身份证}
  \begin{enumerate}
    \item 未经执行机关批准不得离开执行监视居住的处所；（区别取保）
    \item 未经执行机关批准不得会见他人或者通信；（区别取保）
    \item 在传讯的时候及时到案；（同取保）
    \item 不得以任何形式干扰证人作证；（同取保）
    \item 不得毁灭、伪造证据或者串供；（同取保）
    \item 将护照等出入境证件、身份证件、驾驶证件交执行机关保存。（区别取保）
  \end{enumerate}

  \ft{（《刑诉法》第77条）}
\end{frame}

\begin{frame}{【65】严重违反取保候审、监视居住期间义务应当逮捕的情形}{新罪、逃杀、串毁证、还要打击报复他}
  嫌疑人有下列违反取保候审、监视居住规定的行为，检察院应当对嫌疑人予以逮捕：
  \begin{enumerate}
    \item 故意实施新的犯罪；
    \item 企图自杀、逃跑；
    \item 实施毁灭、伪造证据，串供或者干扰证人作证，足以影响侦查、审查起诉工作正常进行；
    \item 对被害人、证人、鉴定人、举报人、控告人及其他人员实施打击报复。
  \end{enumerate}

  \ft{（《高检规则》第101、111条）}
\end{frame}

\begin{frame}{【66】监视居住的“处所”确定}{有家就在家，没家找个家，有家要找家，就得危和恐}
  \begin{enumerate}
    \item 监视居住应当在嫌疑人/被告人的住处执行；无固定住处的，可以在指定的居所执行。
    \item 对于涉嫌危害国家安全犯罪、恐怖活动犯罪，在住处执行可能有碍侦查的，经上一级公安批准，也可以在指定的居所执行。但是，不得在羁押场所、专门的办案场所执行。
  \end{enumerate}

  \ft{（《刑诉法》第75条）}
\end{frame}

\begin{frame}{【67-1】公安可以先行拘留的情形}{正在指认发现证，毁姓逃杀流结多}
  \begin{enumerate}
    \item 正在预备犯罪、实行犯罪或者在犯罪后即时被发觉的；
    \item 被害人或者在场亲眼看见的指认他犯罪的；
    \item 在身边或者住处发现有犯罪证据的；
    \item 有毁灭、伪造证据或者串供可能的；
    \item 不讲真实姓名、住址，身份不明的；
    \item 犯罪后企图自杀、逃跑或者在逃的；
    \item 有流窜作案、多次作案、结伙作案重大嫌疑的。
  \end{enumerate}

  \ft{（《刑诉法》第82条）}
\end{frame}

\begin{frame}{【67-2】检察院在自侦案件的侦查中，以下2种情形可以决定拘留}{自杀、逃跑、在逃的；毁灭、伪造、串供的}
  \begin{enumerate}
    \item 犯罪后企图自杀、逃跑、在逃的；
    \item 有毁灭、伪造证据或者串供可能的。
  \end{enumerate}

  \ft{（《高检规则》第121条）}
\end{frame}

\begin{frame}{【68】拘留后，应当在24小时内完成的程序}{拘留24小时，送看、通知和讯问}
  \begin{enumerate}
    \item 应当立即将被拘留人送看守所羁押，至迟不得超过24小时；
    \item 拘留后24小时以内，通知被拘留人的家属。无法通知、涉嫌危害国家安全犯罪、恐怖活动犯罪通知可能有碍侦查的除外；
    \item 拘留后24小时内，公安、检察院对于被拘留的嫌疑人应当进行讯问。
    \item 【考点提示】谁决定，谁通知，谁决定，谁讯问。
  \end{enumerate}

  \ft{（《刑诉法》第85、86条）}
\end{frame}

\begin{frame}{【69】逮捕后，应当在24小时内完成的程序}{逮捕24小时，仅限通知和讯问}
  应当立即将被逮捕人送看守所羁押；除无法通知的以外，应当在逮捕后24小时以内，通知被逮捕人的家属；法院、检察院对于各自决定逮捕的人，公安对于经检察院批准逮捕的人，都必须在逮捕后的24小时以内进行讯问。

  【考点提示】谁想捕，谁通知，谁想捕，谁讯问。

  \ft{（《刑诉法》第93、94条）}
\end{frame}

\begin{frame}{【70-1】逮捕的3个基本条件}{逮捕需要3条件，有证据有徒刑有危险}
  对有证据证明有犯罪事实，可能判处徒刑以上刑罚的嫌疑人/被告人，采取取保候审尚不足以防止发生社会危险性的，应当予以逮捕。

  \ft{（《刑诉法》第81条）}
\end{frame}

\begin{frame}{【70-2】逮捕的证据条件}{发生了、他干了、已有查证属实了}
  \begin{enumerate}
    \item 有证据证明发生了犯罪事实；（发生了）
    \item 有证据证明该犯罪事实是嫌疑人实施的；（他干了）
    \item 证明其实施犯罪行为的证据已有查证属实的。（已有查证属实了）
  \end{enumerate}

  \ft{（《高检规则》第128条第2、3款）}
\end{frame}

\begin{frame}{【71】批捕前有下列情形之一的，应当讯问嫌疑人}{疑、面、违、难、加认罚、小孩、聋哑、半疯傻}
  \begin{enumerate}
    \item 对是否符合逮捕条件有疑问的；
    \item 嫌疑人要求向检察人员当面陈述的；
    \item 侦查活动可能有重大违法行为的；
    \item 案情重大疑难复杂的；
    \item 嫌疑人认罪认罚的；
    \item 嫌疑人系未成年人的。
    \item 嫌疑人是盲、聋、哑人；
    \item 嫌疑人是尚未完全丧失辨认、控制自己行为能力的精神病人的。
  \end{enumerate}

  \ft{（《高检规则》第280条第1、2款）}
\end{frame}

\begin{frame}{【72】对人大代表适用强制措施的报请程序}{同级直接报，上级要层报，\\下级两路报，两级分别报，\\外地委托报，乡级县里告，\\拘留若已报，逮捕不再报，\\现行先拘后报告}
  \begin{enumerate}
    \item 本级：报请本级人大主席团或者常委会许可；
    \item 上级：层报所属人大同级的检察院报请许可；
    \item 下级：直接报请该代表所属人大主席团或者常委会许可，也可以委托该代表所属的人大同级的检察院报请许可；
    \item 两级：分别依照以上几款的规定报请许可；
    \item 外地：委托该代表所属的人大同级的检察院报请许可；
    \item 乡级：由县级检察院报告乡、民族乡、镇人大；
    \item 拘留已报许可的，逮捕时不需再次报请许可；
    \item 现行犯先行拘留的，拘留后立即依上述途径报告。
  \end{enumerate}

  \ft{（《高检规则》第148、149条）}
\end{frame}

\begin{frame}{【73】逮捕外国人的程序}{危、交、难找最高检，其他先批后备案}
  \begin{enumerate}
    \item 危害国家安全犯罪、涉及国与国之间政治、外交关系以及在适用法律上确有疑难的案件，检察院层报最高检，最高检征求外交部的意见后，作出是否批捕的批复，报送的检察院依据该批复作出是否逮捕的决定。
    \item 其他的涉外案件，应当在作出批准逮捕决定后48小时以内报上一级检察院备案，同时向同级政府外事部门通报。
  \end{enumerate}

  \ft{（《高检规则》第294条）}
\end{frame}

\begin{frame}{【74-1】可以变更逮捕的情形}{疾病、孕乳、唯一扶}
  \begin{enumerate}
    \item 患有严重疾病、生活不能自理的；
    \item 怀孕或者正在哺乳自己婴儿的妇女；
    \item 系生活不能自理的人的唯一扶养人。
  \end{enumerate}

  \ft{（《刑诉解释》169条）}
\end{frame}

\begin{frame}{【74-2】应当变更逮捕的情形}{无罪无责不处罚，量刑非羁押，刑期已折抵，办案已超期}
  \begin{enumerate}
    \item 一审判无罪、不负刑事责任或免予刑事处罚的；
    \item 一审判管制、宣告缓刑、单独适用附加刑，判决尚未生效的；
    \item 被羁押的时间已到一审法院对其判处的刑期期限的；
    \item 案件不能在法律规定的期限内审结的。
  \end{enumerate}

  \ft{（《刑诉解释》170条）}
\end{frame}

\begin{frame}{【75】强制措施的自动解除时间}{单处附加生效后，主刑开始执行后}
  被采取强制措施的被告人，被判处管制、缓刑的，在社区矫正开始后，强制措施自动解除；被单处附加刑的，在判决、裁定发生法律效力后，强制措施自动解除；被判处监禁刑的，在刑罚开始执行后，强制措施自动解除。

  \ft{（《刑诉解释》第172条）}
\end{frame}

\begin{frame}{【76】羁押必要性审查的结案决定主体}{变更“长”决定；不变“官”决定}
  羁押必要性审查的结案决定主体：
  \begin{enumerate}
    \item 经审查认为有继续羁押必要的，由检察官决定结案，并通知办案机关。
    \item 经审查认为无继续羁押必要的，检察官应当报经检察长或者分管副检察长批准，以本院名义向办案机关发出释放或者变更强制措施建议书，并要求办案机关在10日以内回复处理情况。
  \end{enumerate}

  \ft{（《检察院办理羁押必要性审查案件规定》第21、23条）}

  强制措施的继续和变更决定主体：对被告人采取、撤销或者变更强制措施的，由院长决定；决定继续取保候审、监视居住的，可以由合议庭或者独任审判员决定。

  \ft{（《刑诉解释》第147条第2款）}
\end{frame}

\begin{frame}{【77】羁押必要性审查后，检察院应当建议办案机关释放或变更逮捕的情形}{无证据、无徒刑、无必要、已超期}
  \begin{enumerate}
    \item 案件证据发生重大变化，没有证据证明有犯罪事实或犯罪行为系是他所为的；
    \item 案件事实或者情节发生变化，嫌疑人/被告人可能被判处拘役、管制、独立适用附加刑、免予刑事处罚或者判决无罪的；
    \item 案件事实基本查清，证据已收集固定，符合取保候审或者监视居住条件的；
    \item 继续羁押嫌疑人/被告人，羁押期限将超过依法可能判处的刑期的。
  \end{enumerate}

  \ft{（《高检规则》第579条）}
\end{frame}

\begin{frame}{【78】刑事撤回起诉，附带民事诉讼如何处理：}{公诉若撤回，附民调、驳、另；公诉若无罪，附民可调可判可另诉}
  \begin{enumerate}
    \item 刑事撤回起诉，附带民事诉讼如何处理：法院准许检察院撤回起诉的公诉案件，对已经提起的附带民诉，可以进行调解；不宜调解或者经调解不能达成协议的，应当裁定驳回起诉，并告知附带民诉原告人可以另行提起民事诉讼。
    \item 刑事不构成犯罪，附带民事诉讼如何处理：法院认定公诉案件被告人的行为不构成犯罪，对已经提起的附带民诉，经调解不能达成协议的，可以一并作出刑事附带民事判决，也可以告知附带民事原告人另行提起民事诉讼。
  \end{enumerate}

  \ft{（《刑诉解释》第197条）}
\end{frame}

\begin{frame}{【79】人身损害案件中附带民事赔偿范围}{医护交康加误工，残疾辅助死亡葬}
  犯罪行为造成被害人人身损害的，应当赔偿医疗费、护理费、交通费等为治疗和康复支付的合理费用，以及因误工减少的收入。造成被害人残疾的，还应当赔偿残疾生活辅助具费等费用；造成被害人死亡的，还应当赔偿丧葬费等费用。附带民事诉讼当事人就民事赔偿问题达成调解、和解协议的，赔偿范围、数额不受上述限制。

  \ft{（《刑诉解释》第192条）}
\end{frame}

\begin{frame}{【80】附带民事诉前财产保全的特点}{诉前保全有风险，只能申请必须担；\\四八小时作裁定，半月不诉除保全}
  只能向被保全财产所在地、被申请人居住地或者对案件有管辖权的法院申请采取保全措施，法院不能依职权采取。申请人应当提供担保，否则裁定驳回申请。法院接受申请后，必须在48小时内作出裁定。申请人在法院受理刑案后15日内未提起附带民诉的，法院应当解除保全措施。

  \ft{（《刑诉解释》第189条）}
\end{frame}

\begin{frame}{【81】重新计算办案期限的情形}{机关改变；程序重来}
  \begin{enumerate}
    \item 在侦查期间，发现另有重要罪行的，自发现之日起重新计算。
    \item 补充侦查完毕移送检察院或者法院后，重新计算审查起诉期限。
    \item 改变管辖的案件，从改变后的机关收到案件之日起重新计算审理期限。
    \item 二审发回原审的案件，从收到发回的案件之日起，重新计算审理期限。
    \item 简易程序转换为普通程序审理，从决定转为普通程序之日起重新计算。
  \end{enumerate}

  \ft{（《刑诉法》第160、172、175、208、241条《刑诉解释》第368、376条）}
\end{frame}

\begin{frame}{【82】不计入办案期间的特殊的情形}{身份不明、阅卷、中；精神鉴定、监委送}
  \begin{enumerate}
    \item 嫌疑人不讲真实姓名、住址，身份不明的期间。
    \item 二审通知检察院阅卷，1个月阅卷时间不计入。
    \item 中止审理的期间。
    \item 对嫌疑人作精神病鉴定的期间不计入办案期限。
    \item 监委（留置）移送审查起诉的，先行拘留时间不计入审查起诉期限。
  \end{enumerate}

  \ft{（《刑诉法》第160、235、206、149、170条）}
\end{frame}

\begin{frame}{【83】检察院对公安立案监督的步骤}{先要求说理，后通知立案}
  \begin{enumerate}
    \item 要求公安说明不予立案的理由。
    \item 公安在收到通知书的7日内说明不立案的理由。
    \item 认为公安不立案理由不能成立，应当通知公安立案。
    \item 公安在收到《通知立案书》后，应当在15日内决定立案。
  \end{enumerate}

  \ft{（《高检规则》第559-563条）}
\end{frame}

\begin{frame}{【84】讯问嫌疑人的地点}{羁押仅在看守所，未押可在现、住、办}
  除下列情形以外，应当在公安执法办案场所的讯问室进行：
  \begin{enumerate}
    \item 紧急情况下在现场进行讯问的；
    \item 对有严重伤病或者残疾、行动不便的，以及正在怀孕的嫌疑人，在其住处或者就诊的医疗机构进行讯问的。
  \end{enumerate}

  对于已送交看守所羁押的嫌疑人，应当在看守所讯问室进行讯问。

  \ft{（《公安部规定》第198条）}
\end{frame}

\begin{frame}{【85】询问证人的地点}{提、住、侦、现、单}
  侦查人员询问证人，可以在现场进行，也可以到证人所在单位、住处或证人提出的地点进行，在必要的时候，可以通知证人到检察院或者公安提供证言。

  \ft{（《刑诉法》第124条）}
\end{frame}

\begin{frame}{【86-1】公安辨认要求}{7人5物10照片}
  辨认嫌疑人时，被辨认的人数不得少于7人;对嫌疑人照片进行辨认的，不得少于10 人的照片;辨认物品时，混杂的同类物品不得少于5件。

  \ft{（《公安部规定》第260条）}
\end{frame}

\begin{frame}{【86-2】检察院辨认要求}{7人5物10人照5物照}
  辨认嫌疑人时，被辨认的人数不得少于7人，照片不得少于10张；辨认物品时，同类物品不得少于5件，照片不得少于5张。

  \ft{（《高检规则》第226条）}
\end{frame}

\begin{frame}{【87-1】应当邀请见证人的侦查活动}{勘、检、搜、辨、查、扣；必须要有见证人}
  公安开展勘验、检查、搜查、辨认、查封、扣押等侦查活动，应当邀请有关公民作为见证人。

  \ft{（《公安部规定》第194条第1款）}
\end{frame}

\begin{frame}{【87-2】不能担任见证人的情形}{有缺陷、有利害、有职权}
  \begin{enumerate}
    \item 生理上、精神上有缺陷或者年幼，不具有相应辨别能力或者不能正确表达的人；
    \item 与案件有利害关系，可能影响案件公正处理的人；
    \item 公安的工作人员或者其聘用的人员。
  \end{enumerate}

  确因客观原因无法由符合条件的人员担任见证人的，应当对有关侦查活动进行全程录音录像，并在笔录中注明有关情况。

  \ft{（《公安部规定》第194条第2款）}
\end{frame}

\begin{frame}{【88-1】技术侦查的案件范围}{公安：危恐黑毒严；检院：职权侵人权\\监委：贪贿职务犯；都有：追捕在逃犯}
  技术侦查的案件范围：
  \begin{enumerate}
    \item 公安在立案后，对于危害国家安全犯罪、恐怖活动犯罪、黑社会性质的组织犯罪、重大毒品犯罪或者其他严重危害社会的犯罪案件，根据侦查犯罪的需要，经过严格的批准手续，可以采取技术侦查措施。
    \item 检察院在立案后，对于利用职权实施的严重侵犯公民人身权利的重大犯罪案件，根据侦查犯罪的需要，经过严格的批准手续，可以采取技术侦查措施，按照规定交有关机关执行。
    \item 追捕被通缉或者批准、决定逮捕的在逃的嫌疑人/被告人，经过批准，可以采取追捕所必需的技术侦查措施。
  \end{enumerate}

  \ft{（《刑诉法》第150条）}

  技术调查的案件范围：监察机关调查涉嫌重大贪污贿赂等职务犯罪，根据需要，经过严格的批准手续，可以采取技术调查措施，按照规定交有关机关执行。

  \ft{（《监察法》第28条）}
\end{frame}

\begin{frame}{【88-2】技术侦查（调查）的审批有效期}{3 + n个3}
  技术侦查的批准决定自签发之日起3个月以内有效。对于不需要继续的，应当及时解除；对于复杂、疑难案件，经过批准，有效期可以延长，每次不得超过3个月。

  \ft{（《刑诉法》第151条 （《监察法》第28条）}
\end{frame}

\begin{frame}{【88-3】技术侦查的执行程序}{3秘要保密，无关要销毁，不得作他用}
  \begin{enumerate}
    \item 技术侦查措施中知悉的国家秘密、商业秘密和个人隐私，应当保密；无关材料，必须及时销毁。
    \item 技术侦查获取材料，只能用于对犯罪侦查、起诉、审判，不得他用。
  \end{enumerate}

  \ft{（《刑诉法》第152条）}
\end{frame}

\begin{frame}{【89】侦查羁押期限}{（2+1+2+2+x）}
  \begin{enumerate}
    \item 对嫌疑人逮捕后的侦查羁押期限不得超过2个月。
    \item 案情复杂、期限届满不能终结的案件，可以经上一级检察院批准延长1个月。
    \item 下列案件在3个月期限届满不能侦查终结的，经省、自治区、直辖市检察院批准或者决定，可以延长2个月：交通十分不便的边远地区的重大复杂案件；重大的犯罪集团案件；流窜作案的重大复杂案件；犯罪涉及面广，取证困难的重大复杂案件。
    \item 4.对嫌疑人可能判处10年有期徒刑以上刑罚，依照上述规定延长期限届满，仍不能侦查终结的，经省、自治区、直辖市检察院批准或者决定，可以再延长2个月。
    \item 5.因为特殊原因，在较长时间内不宜交付审判的特别重大复杂的案件，由最高检察院报请全国人民代表大会常务委员会批准延期审理。
  \end{enumerate}

  \ft{（《刑诉法》第156-159条）}
\end{frame}

\begin{frame}{【90-1】直接言词原则的特点}{亲自面对、口头陈述}
  是指法官必须在法庭上亲自听取当事人、证人及其他诉讼参与人的口头陈述，案件事实和证据必须由控辩双方当庭口头提出并以口头辩论和质证的方式进行调查。
\end{frame}

\begin{frame}{【90-2】集中审理原则的特点}{人不换、事不断}
  指法院开庭审理案件，应在不更换审判人员的条件下连续进行，不得中断审理的诉讼原则。
\end{frame}

\begin{frame}{【91-1】法定不起诉的情形}{显著轻过时效,特赦告诉和死掉, 还有无罪绝不诉}
  \begin{enumerate}
    \item 嫌疑人实施的行为情节显著轻微，危害不大，不认为是犯罪的；
    \item 嫌疑人的犯罪已过追诉时效期限的；
    \item 嫌疑人的犯罪经特赦令免除刑罚的；
    \item 依照《刑法》告诉才处理的犯罪，没有告诉或者撤回告诉的；
    \item 嫌疑人/被告人死亡的；
    \item 嫌疑人没有犯罪事实。
    \item 其他法律规定免予刑事责任的。
  \end{enumerate}

  \ft{（《刑诉法》第177条第1款 ）}
\end{frame}

\begin{frame}{【91-2】酌定不起诉的条件（同时满足）}{有罪轻微不判刑, 可以酌情不起诉}
  \begin{enumerate}
    \item 嫌疑人的行为已经构成犯罪；
    \item 犯罪情节轻微，依刑法规定不需要判处刑罚或免除刑罚。
  \end{enumerate}

  \ft{（《刑诉法》第177条第2款 ）}
\end{frame}

\begin{frame}{【91-3】存疑不起诉的情形}{证据存疑不起诉，1次补侦可不诉，2次补侦应不诉}
  \begin{enumerate}
    \item 一次补侦后，证据不足，可以作不起诉决定。
    \item 二次补侦后，仍不具备起诉条件，应当不起诉。
  \end{enumerate}

  \ft{（《高检规则》第367条）}
\end{frame}

\begin{frame}{【92】独任审理的适用程序}{基层简速可独任，简易未必是独任，速裁一定要独任}
  \begin{enumerate}
    \item 基层法院适用简易程序、速裁程序的案件可以由审判员一人独任审判。
    \item 适用简易程序审理案件，对可能判处3年有期徒刑以下刑罚的，可以组成合议庭进行审判。
    \item 适用速裁程序的，应当由审判员一人独任审判。
  \end{enumerate}

  \ft{（《刑诉法》第183条第1款、216条第1款、222条第1款 ）}
\end{frame}

\begin{frame}{【93】审判组织的组成方式}{基层简、速可独任；其他统统用合议；\\法官低3高到357；法陪统统3或7；\\最高专业不用陪，地方一审可以陪；\\二审法官3或5，复核法官仅3人；\\再审具体看几审，重审另组合议庭}
  \begin{enumerate}
    \item 基层法院适用简易程序、速裁程序的案件可以由审判员1人独任审判。
    \item 基层、中院第一审案件，由审判员3人或由审判员和陪审员共3人或7人组成合议庭。
    \item 高院审判第一审案件，由审判员3人至7人或由审判员和陪审员共3人或7人组成合议庭。
    \item 最高法审判第一审案件，由审判员3人至7人组成合议庭。
    \item 法院审判上诉和抗诉案件，由审判员3人或5人组成合议庭。
    \item 复核死刑案件，由审判员3人组成合议庭。
    \item 合议庭的成员人数应当是单数。
  \end{enumerate}

  \ft{（《刑诉法》第183、249条）}
\end{frame}

\begin{frame}{【94】下列情形之一的，合议庭成员不承担责任（免责事由）}{无过才无责}
  \begin{enumerate}
    \item 因对法律理解和认识上的偏差而导致案件被改判或者发回重审的；
    \item 因对案件事实和证据认识上的偏差而导致案件被改判或者发回重审的；
    \item 因新的证据而导致案件被改判或者发回重审的；
    \item 因法律修订或者政策调整而导致案件被改判或者发回重审的；
    \item 因裁判所依据的其他法律文书被撤销或变更而导致案件被改判或者发回重审的；
    \item 其他依法履行审判职责不应当承担责任的情形。
  \end{enumerate}

  \ft{（《最高法关于进一步加强合议庭职责的若干规定》第10条）}
\end{frame}

\begin{frame}{【95】合议庭应当提请院长决定提交审委会讨论决定的情形}{抗死要再审，应上审委会}
  \begin{enumerate}
    \item 高、中院拟判处死刑立即执行的案件，以及中院拟判死缓的案件；
    \item 本院已生效的裁判确有错误需要再审的案件；
    \item 检察院依照审判监督程序提出抗诉的案件。
    \item 禁止担任陪审员的情形：
    \item 人民代表大会常务委员会的组成人员，监察委员会、法院、检察院、公安、国家安全机关、司法行政机关的工作人员；
  \end{enumerate}

  \ft{（《刑诉解释》第216条第2款）}
\end{frame}

\begin{frame}{【96】禁止担任陪审员的情形}{公检法司常安监；仲基律师公证员；受刑开除吊执照；纳入失信惩戒严}
  \begin{enumerate}
    \item 律师、公证员、仲裁员、基层法律服务工作者；
    \item 其他因职务原因不适宜担任人民陪审员的人员；
    \item 受过刑事处罚的；
    \item 被开除公职的；
    \item 被吊销律师、公证员执业证书的；
    \item 被纳入失信被执行人名单的；
    \item 因受惩戒被免除人民陪审员职务的；
    \item 其他有严重违法违纪行为，可能影响司法公信的。
  \end{enumerate}

  \ft{（《人民陪审员法》第6、7条）}
\end{frame}

\begin{frame}{【97】应当组成7人合议庭的情形}{征环药十益（甄嬛要十亿）}
  下列第一审案件，由人民陪审员和法官组成7人合议庭进行：
  \begin{enumerate}
    \item 可能判处十年以上有期徒刑、无期徒刑、死刑，社会影响重大的刑事案件；
    \item 根据民事诉讼法、行政诉讼法提起的公益诉讼案件；
    \item 涉及征地拆迁、生态环境保护、食品药品安全，社会影响重大的案件；
    \item 其他社会影响重大的案件。
  \end{enumerate}

  \ft{（《人民陪审员法》第16条）}
\end{frame}

\begin{frame}{【98】公诉案件庭前审查应当退回检察院的情形}{管辖错、人不在、无新证再起诉、16条2-6，通通退回检察院。}
  \begin{enumerate}
    \item 不属本院管辖的；
    \item 被告人不在案的；
    \item 裁定准许撤诉的，没有新的影响定罪量刑的事实、证据，重新起诉的；
    \item 属于刑诉法第16条第2-6项情形的。
  \end{enumerate}

  \ft{（《刑诉解释》第219条）}
\end{frame}

\begin{frame}{【99】法院可以决定召开庭前会议的情形}{多、大、争议}
  \begin{enumerate}
    \item 证据材料较多、案情重大复杂的；
    \item 控辩双方对事实、证据存在较大争议的；
    \item 社会影响重大的；
    \item 需要召开庭前会议的其他情形。
  \end{enumerate}

  \ft{（《刑诉解释》第226条）}
\end{frame}

\begin{frame}{【100】撤诉的理由}{公诉撤诉无罪责，自诉撤诉要自愿}
  \begin{enumerate}
    \item 公诉：检察院要求撤诉的，法院应当审查撤诉理由（无罪、无责），作出是否准许的裁定。
    \item 自诉：自诉人要求撤诉，确实自愿的，应当裁定准许。
  \end{enumerate}

  \ft{（《刑诉解释》第296、329条《高检规则》第424条）}
\end{frame}

\begin{frame}{【101】单位犯罪诉讼代表人的产生}{首选老大，不然其他，最后编外}
  \begin{enumerate}
    \item 被告单位的诉讼代表人，应当是法定代表人、实际控制人或主要负责人。
    \item 上述人员被指控为单位犯罪直接责任人员或因客观原因无法出庭的，应当由被告单位委托其他负责人或者职工作为诉讼代表人。
    \item 难以确定的，可由被告单位委托律师等单位以外的人员担任。
    \item 【考点提示】诉讼代表人不得同时担任被告单位或者被指控为单位犯罪直接责任人员的有关人员的辩护人。
  \end{enumerate}

  \ft{（《刑诉解释》第336条）}
\end{frame}

\begin{frame}{【102】被告单位的诉讼代表人不出庭的，应当按照下列情形分别处理}{老大不来可拘传，其他不来只能换}
  \begin{enumerate}
    \item 诉讼代表人系被告单位的法定代表人、实际控制人或主要负责人，无正当理由拒不出庭的，可以拘传其到庭；因客观原因无法出庭，或者下落不明的，应当要求检察院另行确定诉讼代表人；
    \item 诉讼代表人系其他人员的，应当要求检察院另行确定诉讼代表人出庭。
  \end{enumerate}

  \ft{（《刑诉解释》第337条）}
\end{frame}

\begin{frame}{【103】延期审理的情形}{新证、补侦和回避；法院决定延期审}
  庭审中，遇有下列情形之一，影响审判进行的，可以延期审理：需要通知新的证人到庭，调取新的物证，重新鉴定或者勘验的。检察院要补充侦查，提出建议的。由于当事人申请回避而不能进行审判的。

  \ft{（《刑诉法》第204条）}
\end{frame}

\begin{frame}{【104】中止审理的情形}{他跑、他病、耐心等；法院裁定中止审}
  在审判过程中，有下列情形之一，致使案件在较长时间内无法继续审理的，可以中止审理：被告人患有严重疾病，无法出庭的；被告人脱逃的；自诉人患有严重疾病，无法出庭，未委托诉讼代理人出庭的；由于不能抗拒的原因。

  \ft{（《刑诉法》第206条）}
\end{frame}

\begin{frame}{【105】终止审理的情形}{过时效、特赦、告诉和死掉；无需再等终结审}
  \begin{enumerate}
    \item 嫌疑人的犯罪已过追诉时效期限的；
    \item 经特赦令免除刑罚的；嫌疑人的犯罪已过追诉时效期限的；
    \item 依照刑法告诉才处理的犯罪，没有告诉或者撤回告诉的；嫌疑人的犯罪已过追诉时效期限的；
    \item 嫌疑人/被告人死亡的。
  \end{enumerate}

  \ft{（《刑诉解释》第295条）}
\end{frame}

\begin{frame}{【106】简易程序的适用条件}{清楚、认罪、同意}
  \begin{enumerate}
    \item 基层法院管辖的案件，同时符合下列条件的，法院可以适用简易程序审判：
    \item 案件事实清楚、证据充分的；被告人认罪；被告人对适用简易程序没有异议的。
  \end{enumerate}

  \ft{（《刑诉法》第214条）}
\end{frame}

\begin{frame}{【107】禁止适用简易程序的案件}{盲聋哑、半疯傻、影响大、无罪阿}
  \begin{enumerate}
    \item 被告人是盲、聋、哑人；
    \item 被告人是尚未完全丧失辨认或者控制自己行为能力的精神病人；
    \item 有重大社会影响的；
    \item 共同犯罪案件中部分被告人不认罪或者对适用简易程序有异议的；
    \item 辩护人作无罪辩护的；
    \item 被告人认罪但经审查认为可能不构成犯罪的；
    \item 不宜适用简易程序审理的其他情形。
  \end{enumerate}

  \ft{（《刑诉解释》第360条）}
\end{frame}

\begin{frame}{【108】简易转普通的情形}{无罪、无责、否认、不清}
  \begin{enumerate}
    \item 被告人的行为可能不构成犯罪的；
    \item 被告人可能不负刑事责任的；
    \item 被告人当庭对起诉指控的犯罪事实予以否认的；
    \item 案件事实不清、证据不足的；
    \item 不应当或者不宜适用简易程序的其他情形。
  \end{enumerate}

  \ft{（《刑诉解释》第368条）}
\end{frame}

\begin{frame}{【109】速裁程序的适用条件}{轻微、清楚、认罪、同意}
  基层法院管辖的可能判处3年有期徒刑以下刑罚的案件，案件事实清楚，证据确实、充分，被告人认罪认罚并同意适用速裁程序的，可以适用速裁程序，由审判员一人独任审判。

  \ft{（《刑诉法》第222条）}
\end{frame}

\begin{frame}{【110】速裁程序的禁止范围}{幼聋傻影响大，共犯异议民未达，无罪辩护或其他}
  \begin{enumerate}
    \item 被告人是未成年人；
    \item 被告人是盲、聋、哑人；
    \item 被告人是尚未完全丧失辨认控制自己行为能力的精神病人；
    \item 案件有重大社会影响；
    \item 共犯中部分被告人对指控的事实、罪名、量刑建议或适用速裁程序有异议；
    \item 被告人与被害人或其法定代理人没有就附带民事诉讼赔偿等事项达成调解或和解协议；
    \item 辩护人作无罪辩护的；
    \item 其他不宜适用速裁程序的情形。
  \end{enumerate}

  \ft{（《刑诉解释》第370条）}
\end{frame}

\begin{frame}{【111】速裁程序应当转为普通或简易程序审理的情形}{无罪无责违意愿、否认疑难争议大}
  \begin{enumerate}
    \item 被告人的行为可能不构成犯罪或者不应当追究刑事责任的；
    \item 被告人违背意愿认罪认罚的；
    \item 被告人否认指控的犯罪事实的；
    \item 案件疑难、复杂或者对适用法律有重大争议的；
    \item 其他不宜适用速裁程序的情形。
  \end{enumerate}

  \ft{（《刑诉解释》第375条）}
\end{frame}

\begin{frame}{【112】判决、裁定、决定的区别}{裁定好比墙头草，实体程序都能搞；\\公检只能做决定，判决裁定法院行；\\判决只能书面出，生效判决要唯一；\\决定做出就生效，不能上诉不能抗}
  \begin{enumerate}
    \item 判决只能对实体，决定只能对程序，裁定既能对实体又能对程序。
    \item 判决、裁定只能由法院做出，决定可以由公、检、法三机关做出，但公、检只能做决定。
    \item 判决必须是书面形式，裁定、决定可以是口头形式。
    \item 生效的判决要求唯一，裁定、决定可以为多个。
    \item 判决、裁定可以上诉、抗诉，决定做出就生效，不服只能申请复议。
  \end{enumerate}
\end{frame}

\begin{frame}{【113】上诉和抗诉的差异}{上诉灵活抗诉严，无理口头两级院，有理书面原法院}
  \begin{enumerate}
    \item 上诉：无需理由，书面口头都可以，上诉状既可以向原法院也可以向上一级法院提交。
    \item 抗诉：需要理由，只能用书面方式，只能通过原法院提交抗诉书。
  \end{enumerate}

  \ft{（《刑诉法》第227-232条）}
\end{frame}

\begin{frame}{【114】上诉不加刑原则的适用对象}{被上被抗可加刑；其他通通不加刑。}
  第二审法院审理被告人或者他的法定代理人、辩护人、近亲属上诉的案件，不得加重被告人的刑罚。检察院提出抗诉或者自诉人提出上诉的，不受前款规定的限制。

  \ft{（《刑诉法》第237条、《刑诉解释》第402条）}
\end{frame}

\begin{frame}{【115-1】一审后引发的二审必须开庭审理的情形}{抗、死、事证必开庭}
  \begin{enumerate}
    \item 一审上诉案件对认定的事实、证据提出异议，可能影响定罪量刑的；
    \item 被告人被判处死刑的上诉案件；（含死缓）
    \item 检察院抗诉的案件；
  \end{enumerate}

  \ft{（《刑诉解释》第393条）}
\end{frame}

\begin{frame}{【115-2】发回重审的二审必须开庭的情形}{事实证据、纠程序}
  最高法不予核准发回第二审法院重新审判的案件，第二审法院可以直接改判；必须通过开庭查清事实、核实证据或者纠正原审程序违法的，应当开庭审理。

  \ft{（《刑诉解释》第418条）}
\end{frame}

\begin{frame}{【116】二审的裁判结果}{二审结局有三种：维持、改判和发回}
  \begin{enumerate}
    \item 维持原判：事实清楚，证据充分，适用法律正确，量刑适当。
    \item 只改不发：原判决认定事实没有错误，但适用法律有错误或者量刑不当的。
    \item 可改可发：原判决事实不清楚或者证据不足的。（限发1次）
    \item 只发不改：违反法律规定的诉讼程序的。（无限次）
  \end{enumerate}

  \ft{（《刑诉法》第236、238条）}
\end{frame}

\begin{frame}{【117】法定刑以下量刑或死刑复核案件的上诉、上报路线}{上诉抗诉走二审；不上不抗就上报}
  \begin{enumerate}
    \item 一审后上诉、抗诉的，应当按二审程序审理。二审如果维持原判，再逐级报最高法核准。
    \item 一审后没有提起上诉、抗诉的，期满将逐级报请至最高法核准。（层层上报）
  \end{enumerate}

  \ft{（《刑诉解释》第414、423条）}
\end{frame}

\begin{frame}{【118】复核案件在上报途中的处理问题}{层层上报，上报路上不改判}
  \begin{enumerate}
    \item 上级法院同意原判的，应当逐级报请最高法核准；（层层上报）
    \item 上级法院不同意原判的，应当裁定发回重新审判或按照第二审程序提审。
    \item 【考点提示】上报途中任何一个上级法院都有否决权，但不能直接改判，要么发回重审，要么提审。（上报路上不改判）
  \end{enumerate}

  \ft{（《刑诉解释》第414、423条）}
\end{frame}

\begin{frame}{【119】最高法核准死刑的结果}{没错核准，瑕疵纠正，见错发回，不死必要时可改}
  \begin{enumerate}
    \item 没错核准：原判事实、法律正确、量刑适当、诉讼程序合法的。
    \item 瑕疵纠正：原判死刑并无不当，但具体认定某一事实或引用的法律条款等不完全准确、规范的，可以在纠正后核准死刑。
    \item 见错发回：
          \begin{enumerate}
            \item 原判事实不清、证据不足的；（事实证据问题）
            \item 复核期间出现新的影响定罪量刑的事实、证据的；（事实证据问题）
            \item 原判认定事实正确，但依法不应当判处死刑的；（量刑问题）
            \item 原审违反法定程序，可能影响正确判决的。（程序问题）
          \end{enumerate}
    \item 不死必要时可改：原判认定事实正确，但依法不应当判处死刑的，根据案件情况，必要时，也可以依法改判。
  \end{enumerate}

  \ft{（《刑诉解释》第429条）}
\end{frame}

\begin{frame}{【120】死缓复核中对量刑错误的处理}{变轻容易变重难}
  \begin{enumerate}
    \item 变轻容易：高院复核死缓案件，认为量刑过重，应当改判。
    \item 变重难：高院复核死缓案件，认为量刑过轻，不能直接改判，确有必要改判的，只能在核准死缓后，按照审判监督程序重新审判。（复核死刑缓期执行案件，不得加重被告人的刑罚）
  \end{enumerate}

  【考点提示】此考点和二审的处理有细微差别，二审中如果被上、被抗还是可以加刑的，但死缓复核程序中一律不能加重刑罚。

  \ft{（《刑诉解释》第428条）}
\end{frame}

\begin{frame}{【121】再审程序应当开庭的情形}{一、抗、事证、加重罚}
  \begin{enumerate}
    \item 依照第一审程序审理的。
    \item 依照第二审程序需要对事实或者证据进行审理的。
    \item 检察院按照审判监督程序提出抗诉的。
    \item 可能对原审被告人（原审上诉人）加重刑罚的。
    \item 有其他应当开庭审理情形的。
  \end{enumerate}

  \ft{（《最高法关于刑事再审案件开庭审理程序的具体规定》第5条）}
\end{frame}

\begin{frame}{【122】法院执行的刑种}{收钱、杀人、不干活}
  死刑立即执行、罚金、没收财产、及无罪、免除刑罚。

  \ft{（《刑诉法》第260条《刑诉解释》第499、522条）}
\end{frame}

\begin{frame}{【123】监狱执行的刑种}{徒刑监禁3月起}
  死刑缓期二年执行、无期徒刑、有期徒刑（余刑超3月以上）。

  \ft{（《刑诉法》第264条）}
\end{frame}

\begin{frame}{【124】公安执行的刑种}{不足3、拘、夺政权}
  余刑不足3个月的有期徒刑和拘役、剥夺政治权利。

  \ft{（《刑诉法》第264、270条）}
\end{frame}

\begin{frame}{【125】社区矫正机构执行的刑种}{社区最闲陪你玩}
  管制、宣告缓刑、假释或暂予监外执行。

  \ft{（《刑诉法》第269条）}
\end{frame}

\begin{frame}{【126】财产不足以支付的，按下列顺序执行}{医、赔、债、罚、没}
  \begin{enumerate}
    \item 人身损害赔偿中的医疗费用；
    \item 退赔被害人的损失；
    \item 其他民事债务；
    \item 罚金；
    \item 没收财产。
  \end{enumerate}

  债权人对执行标的依法享有优先受偿权，其主张优先受偿的，法院应当在前款第项规定的医疗费用受偿后，予以支持。

  \ft{（《关于刑事裁判涉财产部分执行的若干规定》第13条）}
\end{frame}

\begin{frame}{【127】下列情形之一，应当暂停执行死刑，并层报最高法院}{其他、共犯、立功、孕}
  \begin{enumerate}
    \item 罪犯可能有其他犯罪的；
    \item 共犯到案，可能影响罪犯量刑的；
    \item 共犯暂停，可能影响罪犯量刑的；
    \item 重大立功，可能需要改判的；
    \item 罪犯怀孕的；
    \item 判决、裁定可能有影响定罪量刑的其他错误的。
  \end{enumerate}

  \ft{（《刑诉解释》第500条）}
\end{frame}

\begin{frame}{【128】监外执行的条件}{疾病、孕乳、不自理}
  \begin{enumerate}
    \item 有严重疾病需要保外就医的。
    \item 怀孕或者正在哺乳自己婴儿的妇女。
    \item 生活不能自理，适用暂予监外执行不致危害社会的。
  \end{enumerate}

  \ft{（《刑诉法》第265条）}
\end{frame}

\begin{frame}{【129】下列减刑、假释案件，应当开庭审理}{重大立功不一般；影响重大有意见；有权有势又有钱}
  \begin{enumerate}
    \item 因罪犯有重大立功表现提请减刑的；
    \item 提请减刑的起始时间、间隔时间或减刑幅度不符合一般规定的；
    \item 罪犯系职务犯罪罪犯，组织、领导、参加、包庇、纵容黑社会性质组织罪犯，破坏金融管理秩序罪犯或者金融诈骗罪犯的；
    \item 社会影响重大或者社会关注度高的；
    \item 公示期间收到不同意见的；
    \item 检察院提出异议的；
    \item 有必要开庭审理的其他案件。
  \end{enumerate}

  \ft{（《刑诉解释》第538条）}
\end{frame}

\begin{frame}{【130】附条件不起诉的条件（同时满足）}{小孩、犯罪、456； 1年以下、够起诉；悔过、可以附条件}
  \begin{enumerate}
    \item 对于犯罪时已满14周岁不满18周岁的未成年人；
    \item 涉嫌刑法分则第四章、第五章、第六章规定的犯罪；
    \item 根据具体犯罪事实、情节，可能被判处1年有期徒刑以下刑罚；
    \item 犯罪事实清楚，证据确实、充分，符合起诉条件；
    \item 具有悔罪表现；
    \item 检察院可以作出附条件不起诉决定。
  \end{enumerate}

  \ft{（《刑诉法》第282条）}
\end{frame}

\begin{frame}{【131】被附条件不起诉的未成年嫌疑人在考验期应当遵守下列规定}{市县、报告、守法、教育}
  \begin{enumerate}
    \item 遵守法律法规，服从监督；
    \item 按照考察机关的规定报告自己的活动情况；
    \item 离开所居住的市、县或者迁居，应当报经考察机关批准；
    \item 按照考察机关的要求接受矫治和教育。
  \end{enumerate}

  \ft{（《高检规则》第475条）}
\end{frame}

\begin{frame}{【132】附条件不起诉考验后应当起诉的情形}{新罪、漏罪、违法、违规}
  \begin{enumerate}
    \item 在考验期内有下列情形之一的，检察院应当撤销附条件不起诉的决定，提起公诉：
          \begin{enumerate}
            \item 实施新的犯罪或者发现决定附条件不起诉以前还有其他犯罪需要追诉的；
            \item 违反治安管理规定或者考察机关有关附条件不起诉的监督管理规定，情节严重的。
          \end{enumerate}
    \item 在考验期内没有上述情形，考验期满的，检察院应当作出不起诉的决定。
  \end{enumerate}

  \ft{（《刑诉法》第284条）}
\end{frame}

\begin{frame}{【133】公诉案件可以可以和解的案件范围}{民间人财3年下，7年过失渎职外}
  因民间纠纷引起，涉嫌刑法分则第4章、第5章规定的犯罪案件，可能判处3年有期徒刑以下刑罚的；除渎职犯罪以外的可能判处7年有期徒刑以下刑罚的过失犯罪案件。

  \ft{（《刑诉法》第288条）}
\end{frame}

\begin{frame}{【134】没收违法所得程序的适用条件}{贪恐重大逃一年，死了直接收赃钱}
  对于贪污贿赂犯罪、恐怖活动犯罪等重大犯罪案件，嫌疑人/被告人逃匿，在通缉1年后不能到案，或者嫌疑人/被告人死亡，依照刑法规定应当追缴其违法所得及其他涉案财产的，检察院可以向法院提出没收违法所得的申请。

  \ft{（《刑诉法》第298条）}
\end{frame}

\begin{frame}{【135】检察院对公安应当启动而未启动没收或强疗程序的监督}{先要求说理，后通知启动}
  \begin{enumerate}
    \item 公安应当启动而不启动，可以要求公安7日以内书面说明不启动的理由。
    \item 经审查，认为公安不启动理由不能成立的，应当通知公安启动程序。
  \end{enumerate}

  \ft{（《高检规则》第524、540条）}
\end{frame}

\begin{frame}{【136】强疗程序的适用条件}{暴力、无责、有危险}
  实施暴力行为，危害公共安全或者严重危害公民人身安全，经法定程序鉴定依法不负刑事责任的精神病人，有继续危害社会可能的，可以予以强制医疗。

  \ft{（《刑诉法》第302条）}
\end{frame}

\begin{frame}{【137】认罪认罚的适用范围}{认罪认罚无禁区}
  \begin{enumerate}
    \item 贯穿刑诉全过程，适用于侦查、起诉、审判各个阶段。
    \item 认罪认罚制度没有适用罪名和可能判处刑罚的限定。
  \end{enumerate}

  \ft{（《关于适用认罪认罚从宽制度的指导意见》第5条）}
\end{frame}

\begin{frame}{【138】不需要签署认罪认罚具结书的情形}{盲聋哑、半疯傻、小孩法辩不认罚}
  嫌疑人自愿认罪，同意量刑建议和程序适用的，应当在辩护人或者值班律师在场的情况下签署认罪认罚具结书。嫌疑人认罪认罚，有下列情形之一的，不需要签署认罪认罚具结书：
  \begin{enumerate}
    \item 嫌疑人是盲、聋、哑人，或者是尚未完全丧失辨认或者控制自己行为能力的精神病人的；
    \item 未成年嫌疑人的法定代理人、辩护人对未成年人认罪认罚有异议的；其他不需要签署认罪认罚具结书的情形。
  \end{enumerate}

  \ft{（《刑诉法》第174条）}
\end{frame}

\begin{frame}{【139】检察院对认罪认罚案件自愿性、合法性的审查}{愿常理、告明诚；违意愿、可重来}
  侦查阶段认罪认罚的案件，审查起诉阶段检察院应当重点审查以下内容：是否自愿；是否正常;是否理解；是否告知；是否写明；是否真诚。经审查，嫌疑人违背意愿认罪认罚的，检察院可以重新开展认罪认罚工作。

  \ft{（《认罪认罚指导意见》第28条）}
\end{frame}

\begin{frame}{【140-1】认罪认罚法院对检察院指控的罪名和量刑建议不采纳的情形}{无罪无责违意愿，否认不一或其他，罪刑建议不采纳}
  对于认罪认罚案件，法院依法作出判决时，一般应当采纳检察院指控的罪名和量刑建议，但有下列情形的除外：
  \begin{enumerate}
    \item 被告人的行为不构成犯罪或者不应当追究其刑事责任的；
    \item 被告人违背意愿认罪认罚的；
    \item 被告人否认指控的犯罪事实的；
    \item 起诉指控的罪名与审理认定的罪名不一致的；
    \item 其他可能影响公正审判的情形。
  \end{enumerate}

  \ft{（《刑诉法》第201条第1款）}
\end{frame}

\begin{frame}{【140-2】认罪认罚案件中检察院可以调整量刑建议的情形}{量刑建议显不当，应当建议其调整，仍然不当依法判}
  法院经审理认为量刑建议明显不当，或者被告人、辩护人对量刑建议提出异议的，法院应当告知检察院，检察院可以调整量刑建议。
  检察院不调整量刑建议或者调整量刑建议后仍然明显不当的，法院应当依法作出判决。

  \ft{（《刑诉法》第201条、《关于适用认罪认罚从宽制度的指导意见》第41条）}
\end{frame}

\begin{frame}{【141】认罪认罚之后反悔的处理}{认罪之后若反悔，一夜回到认罪前}
  \begin{enumerate}
    \item 提起公诉前反悔的：具结书失效，检察院应当在全面审查事实证据基础上，提起公诉。
    \item 酌定不起诉后反悔：
          \begin{enumerate}
            \item 确实没有犯罪事实，或符合刑诉法第16条情形的，应当撤销原（酌定）不起诉决定，依照重新作出（法定）不起诉决定；
            \item 犯罪情节轻微，依刑法不需要判处刑罚或免除刑罚的，可以维持原（酌定）不起诉决定；
            \item 排除认罪认罚因素后，符合起诉条件的，应当撤销原不起诉决定，依法提起公诉。
          \end{enumerate}
    \item 审判程序中反悔的：法院应当根据审理查明的事实，依法作出裁判。如果是简易、速裁程序则需要转换为普通程序。
  \end{enumerate}

  \ft{（《认罪认罚指导意见》第51、52、53条）}
\end{frame}

\begin{frame}{【142】认罪认罚案件可以免于刑事处罚的条件}{如实供、大利功、高检核}
  嫌疑人自愿如实供述涉嫌犯罪的事实，有重大立功或者案件涉及国家重大利益的，经最高检察院核准，公安可以撤销案件，检察院可以作出不起诉决定，也可以对涉嫌数罪中的一项或者多项不起诉。

  \ft{（《刑诉法》第182条）}
\end{frame}

\begin{frame}{【143-1】潜逃境外案件缺席审判的条件}{贪贿只要在境外；危恐境外及时审，还要最高检核准。}
  \begin{enumerate}
    \item 贪污贿赂犯罪，嫌疑人/被告人潜逃境外。
    \item 严重的危害国家安全、恐怖活动犯罪，嫌疑人/被告人潜逃境外，需要及时审判，且经过最高检核准的案件。
  \end{enumerate}

  \ft{（《刑诉法》第291条第1款）}
\end{frame}

\begin{frame}{【143-2】因被告人严重疾病缺席审判的条件}{疾病中止超6月，还要申请或同意。}
  因被告人患有严重疾病无法出庭，中止审理超过6个月，被告人仍无法出庭，被告人及其法定代理人、近亲属申请或者同意恢复审理的，法院可以在被告人不出庭的情况下缺席审理。

  \ft{（《刑诉法》第296条）}
\end{frame}

\begin{frame}{【143-3】因被告人死亡缺席审判的条件}{被告死亡若缺席，除非因为他无罪。若是审监他死亡，有罪也能缺席判。}
  \begin{enumerate}
    \item 被告人死亡的，法院应当裁定终止审理，但有证据证明被告人无罪，法院经缺席审理确认无罪的，应当依法作出判决。
    \item 法院按照审判监督程序重新审判的案件，被告人死亡的，法院可以缺席审理，依法作出判决。
  \end{enumerate}

  \ft{（《刑诉法》第297条）}
\end{frame}
\end{document}
